\chapter{Monge--Amp\`ere energy and torsion}
\label{chap:energy}

In the K\"ahler setting, the Monge--Amp\`ere operator is the differential of a scalar energy. This relation is essential in variational arguments and in the geometry of the space of potentials. Before attempting to extend those arguments, we should test the necessary symmetry of the linearization.

We work on the open convex set $\HH_\omega$ of smooth strictly positive potentials. A tangent direction is a smooth real function. Define the one-form
\[
 \alpha_u(v)=\int_Xv\,\omega_u^n.
\]
A primitive of $\alpha$ means a functional $E$ whose derivative in every smooth direction is $DE_u(v)=\alpha_u(v)$. In discussing second derivatives it suffices to require that the restriction of $E$ to every finite-dimensional affine plane be twice continuously differentiable.

\section{The symmetry test}

\begin{proposition}\label{prop:energy-symmetry}
If $\alpha$ has a primitive near $u\in\HH_\omega$, then for every smooth $v,w$,
\begin{equation}\label{eq:energy-symmetry}
 \int_Xw\ddc v\wedge\omega_u^{n-1}
 =\int_Xv\ddc w\wedge\omega_u^{n-1}.
\end{equation}
In particular, $\ddc(\omega_u^{n-1})=0$.
\end{proposition}
\begin{proof}
For small $s,t$, consider the twice differentiable function $E(u+sv+tw)$. Differentiating its asserted first derivatives gives
\[
 \partial_s\partial_tE(u+sv+tw)\big|_{s=t=0}
 =n\int_Xw\ddc v\wedge\omega_u^{n-1},
\]
and the opposite order gives the right side of \eqref{eq:energy-symmetry}, multiplied by $n$. Equality of mixed derivatives proves the identity.

Set $w=1$. The right side vanishes. By \cref{lem:ibp}, the left side is $\int_Xv\ddc(\omega_u^{n-1})$. Its vanishing for every $v$ proves the last assertion.
\end{proof}

\begin{corollary}\label{cor:no-primitive}
If $\ddc(\omega^{n-1})\neq0$, the raw Monge--Amp\`ere one-form has no primitive on a neighborhood of $0$ in $\HH_\omega$. In particular, this applies to the background of \cref{ex:torus-mass}.
\end{corollary}
\begin{proof}
Apply \cref{prop:energy-symmetry} at $u=0$. In the torus example the first variation of mass is nonzero, so \cref{prop:mass-variation} gives $\ddc\omega\neq0$.
\end{proof}

Vanishing of the top-degree form $\ddc(\omega^{n-1})$ tests only the constant tangent direction. Full symmetry in \eqref{eq:energy-symmetry} is a stronger requirement. We now show that requiring a primitive on an open set forces closedness of the background itself.

\begin{lemma}\label{lem:formal-symmetry}
For a smooth real $(n-1,n-1)$-form $\Omega$, the pairing
$B(v,w)=\int_Xw\ddc v\wedge\Omega$ is symmetric on smooth real functions if and only if $\mathrm{d}\Omega=0$.
\end{lemma}
\begin{proof}
Integration by parts shows that symmetry is equivalent to
\[
 \ddc(w\Omega)-\ddc w\wedge\Omega
 =w\ddc\Omega+\mathrm{d}w\wedge\dc\Omega-\dc w\wedge \mathrm{d}\Omega=0
\]
for every smooth $w$. Taking $w=1$ gives $\ddc\Omega=0$. The remaining expression is linear in the first jet of $w$, so we can complexify the test functions. At a chosen point take a cutoff holomorphic linear function with $\partial w=\lambda$ arbitrary and $\bar\partial w=0$. With $a=\mathrm i/(4\pi)$, the first-order expression is
\[
 2a\bigl(\partial w\wedge\bar\partial\Omega
             -\bar\partial w\wedge\partial\Omega\bigr).
\]
It follows that $\lambda\wedge\bar\partial\Omega=0$ for every $(1,0)$-covector $\lambda$. The pairing with $(n-1,n)$-forms is nondegenerate, hence $\bar\partial\Omega=0$. Reality gives $\partial\Omega=0$ as well. Conversely, if $\mathrm{d}\Omega=0$, its pure type implies $\dc\Omega=0$, and the same product formula proves symmetry.
\end{proof}

\begin{lemma}\label{lem:torsion-polarization}
Let $T$ be a form of type $(2,1)+(1,2)$ on a complex vector space of dimension $n\geq2$. If $T\wedge A^{n-2}=0$ for every real $(1,1)$-form $A$ in an open neighborhood of a positive form, then $T=0$.
\end{lemma}
\begin{proof}
For $n=2$ the assertion is immediate. Otherwise the expression is polynomial in the real coefficients of $A$, so it vanishes for every real $(1,1)$-form. Polarization and then complexification give
$T\wedge A_1\wedge\cdots\wedge A_{n-2}=0$ for arbitrary complex $(1,1)$-forms $A_j$. Such products span all forms of bidegree $(n-2,n-2)$.

The two type components of $T$ vanish separately after these wedges. Wedge the identity for $T^{2,1}$ with arbitrary $(0,1)$-covectors. The resulting test forms span bidegree $(n-2,n-1)$, which pairs perfectly with bidegree $(2,1)$ to top degree. Thus $T^{2,1}=0$. The other component is treated with $(1,0)$-covectors.
\end{proof}

\begin{theorem}\label{thm:raw-primitive}
For $n\geq2$, the raw Monge--Amp\`ere one-form has a primitive on $\HH_\omega$ if and only if $\mathrm{d}\omega=0$. In fact, a primitive on a $C^2$-open neighborhood of $0$ already forces $\mathrm{d}\omega=0$.
\end{theorem}
\begin{proof}
Suppose a primitive exists on such a neighborhood $U$. By \cref{prop:energy-symmetry,lem:formal-symmetry}, for every $u\in U$,
\[
 0=\mathrm{d}(\omega_u^{n-1})
 =(n-1)\mathrm{d}\omega\wedge(\omega+\ddc u)^{n-2}.
\]
Fix a point $x$. A real quadratic polynomial in local coordinates, multiplied by a cutoff equal to one near $x$, realizes any prescribed small real $(1,1)$-form $H=\ddc u(x)$. These functions can be chosen linearly in $H$; their global $C^2$ norms are bounded by a fixed constant times the coefficient norm of $H$. Thus sufficiently small $H$ are realized by potentials in $U$. The preceding identity holds with $\omega(x)+H$ ranging over an open set. Apply \cref{lem:torsion-polarization} with $T=\mathrm{d}\omega(x)$ to conclude that $\mathrm{d}\omega(x)=0$.

Since $x$ was arbitrary, $\mathrm{d}\omega=0$. Conversely, \cref{prop:closed-energy} below constructs a primitive for a closed background.
\end{proof}

At a single potential the symmetry criterion is $\mathrm{d}(\omega_u^{n-1})=0$, the balanced condition. Requiring it throughout an open neighborhood is stronger and forces $\mathrm{d}\omega=0$. In particular, the Gauduchon condition $\ddc(\omega^{n-1})=0$ does not remove the variational obstruction.

There are constructions of torsion-corrected energy functionals in the Hermitian literature, including work of Yi Li \cite[Equation~(1.18) and Theorem~1.2]{YiLi10}. Such constructions do not contradict this obstruction: a correction changes the one-form being integrated. To use a corrected energy in a variational problem, one must compute its derivative and verify that its critical-point equation is the desired equation.

\section{The closed-background calculation}

The standard primitive makes the role of closedness visible. For $\mathrm{d}\omega=0$, define
\begin{equation}\label{eq:closed-energy}
 E_\omega(u)=\frac1{n+1}\sum_{j=0}^n
 \int_Xu\,\omega_u^j\wedge\omega^{n-j}.
\end{equation}

\begin{proposition}\label{prop:closed-energy}
If $\mathrm{d}\omega=0$, then $DE_{\omega,u}(v)=\int_Xv\omega_u^n$.
\end{proposition}
\begin{proof}
Differentiate the finite sum in \eqref{eq:closed-energy}. The term indexed by $j$ contributes
\[
 \int_Xv\omega_u^j\wedge\omega^{n-j}
 +j\int_Xu\ddc v\wedge\omega_u^{j-1}\wedge\omega^{n-j},
\]
with the second summand absent for $j=0$. For $j\geq1$, the form $\omega_u^{j-1}\wedge\omega^{n-j}$ is closed. By \cref{lem:formal-symmetry}, $u$ and $v$ may therefore be interchanged in the second summand. Using $\ddc u=\omega_u-\omega$, the contribution becomes
\[
 (j+1)\int_Xv\omega_u^j\wedge\omega^{n-j}
 -j\int_Xv\omega_u^{j-1}\wedge\omega^{n-j+1}.
\]
The sum telescopes to $(n+1)\int_Xv\omega_u^n$. Division by $n+1$ completes the proof.
\end{proof}

This calculation uses closedness of the background, not the existence of any additional K\"ahler metric. The positivity needed to regard the domain as a space of positive potentials is a separate condition.

\section{Integrating along radial paths}

Even when a one-form has no primitive, it can be integrated along a chosen family of paths. Since $\HH_\omega$ is convex and contains $0$, define
\[
 F_\omega(u)=\int_0^1\alpha_{tu}(u)\,\mathrm{d}t
 =\int_0^1\int_Xu(\omega+t\ddc u)^n\,\mathrm{d}t.
\]
This is a well-defined polynomial functional of a smooth potential. Its existence does not imply $DF_\omega=\alpha$.

\begin{proposition}\label{prop:radial-surface}
In complex dimension two,
\[
 F_\omega(u)=\int_Xu\omega^2
 +\int_Xu\omega\wedge\ddc u
 +\frac13\int_Xu(\ddc u)^2.
\]
For a smooth direction $v$, its differential satisfies
\begin{equation}\label{eq:radial-defect}
 DF_{\omega,u}(v)-\int_Xv\omega_u^2
 =\int_Xv\bigl(u\ddc\omega+\mathrm{d}u\wedge\dc\omega
                      -\dc u\wedge \mathrm{d}\omega\bigr).
\end{equation}
\end{proposition}
\begin{proof}
The formula for $F_\omega$ follows by expanding the square and integrating $1,t,t^2$ on $[0,1]$.

Its first term differentiates to $\int_Xv\omega^2$. The second differentiates to
\[
 \int_Xv\omega\wedge\ddc u+\int_Xu\omega\wedge\ddc v.
\]
The derivative of the last term is $\int_Xv(\ddc u)^2$: indeed, the direct derivative gives one copy with coefficient $1/3$, while integration by parts against the closed form $\ddc u$ changes each of the other two copies into the same integral.

Subtracting the expansion of $\int_Xv\omega_u^2$ leaves
\[
 \int_Xu\omega\wedge\ddc v-\int_Xv\omega\wedge\ddc u.
\]
By \cref{lem:ibp}, the first integral equals $\int_Xv\ddc(u\omega)$. Expanding $\ddc(u\omega)$ gives \eqref{eq:radial-defect}:
\[
 \ddc(u\omega)
 =\ddc u\wedge\omega-\dc u\wedge \mathrm{d}\omega
    +\mathrm{d}u\wedge\dc\omega+u\ddc\omega.
\]
\end{proof}

Subtracting a linear measure term from $F_\omega$ gives an Euler equation containing the additional expression in \eqref{eq:radial-defect}. A corrected variational construction must cancel this expression in its first variation.

\section{Consequences for a theory of singularities}

The source book defines a distance on singularity types using the geometry of geodesic rays. A Hermitian construction must distinguish several tasks: defining a suitable class of subgeodesics, constructing envelopes, establishing an energy with a specified derivative, and proving metric identities. A formal positivity condition on a product does not accomplish all four tasks.

The calculations above give two tests for any proposed extension. First, restricting to smooth potentials must reproduce the torsion terms. Second, a claimed primitive of the raw operator must satisfy \eqref{eq:energy-symmetry}. These tests are useful before passing to unbounded potentials, where approximation can hide the same obstruction.

The chapters on interpolation and test curves develop the order and envelope constructions. The energy and metric identities require further arguments; their precise status is recorded there.
